\documentclass[12pt,a4paper]{article}

% Modelo fornecido para o trabalho: manter este arquivo na mesma pasta.
% ===========================================================================
% QUICK EXAMPLES
% ===========================================================================
%
% The examples below show how to use the specification template.
% They can be copied and adapted directly into the main document.
%
% ---------------------------------------------------------------------------
% Functional Requirement
% ---------------------------------------------------------------------------
%
% \fr{User Authentication}{
%     \item The system shall allow users to authenticate using
%     their username and password.
%
%     \item The system shall reject invalid credentials.
%
%     \item The system shall lock the account after three
%     consecutive failed attempts.
% }
%
% ---------------------------------------------------------------------------
% Non-functional Requirement
% ---------------------------------------------------------------------------
%
% \nfr{Authentication Response Time}{
%     \item The system shall authenticate users within 2 seconds
%     under normal operating conditions.
% }
%
% ---------------------------------------------------------------------------
% Use Case
% ---------------------------------------------------------------------------
%
% \begin{usecase}{Perform Login}
%
%     \ucdescription{
%         Allows a user to authenticate into the system using
%         their credentials.
%     }
%
%     \ucactors{
%         \item User
%     }
%
%     \ucpreconditions{
%         \item The user must be registered in the system.
%         \item The system must be available.
%     }
%
%     \ucmainflow{
%         \item The user enters their username and password.
%         \item The system validates the provided credentials.
%         \item The system creates an authenticated session.
%         \item The system grants access to the application.
%     }
%
%     \ucalternativeflow{
%         \item If the credentials are invalid, the system informs
%         the user that the username or password is incorrect.
%
%         \item If the account is locked, the system denies access.
%     }
%
%     \ucbusinessrules{
%         \item The provided password must match the registered password.
%         \item The account is locked after three consecutive failed attempts.
%     }
%
%     \ucpostconditions{
%         \item The user is authenticated.
%         \item An authenticated session exists for the user.
%     }

%     \ucrelatedrequirements{
%         \item FR01
%     }
%
% \end{usecase}
%
% ---------------------------------------------------------------------------
% Operation Contract
% ---------------------------------------------------------------------------
%
% \begin{operationcontract}{Create Order}
%
%     \ocoperation{
%         createOrder(customer)
%     }
%
%     \ocdescription{
%         Creates a new order for the specified customer.
%     }
%
%     \ocpreconditions{
%         \item The customer must exist.
%         \item The customer must be active.
%     }
%
%     \ocpostconditions{
%         \item A new order is created.
%         \item The order is associated with the customer.
%         \item The order status is set to \texttt{OPEN}.
%     }
%
%     \ocrelatedusecases{
%         \item \ref{uc:place-order}
%     }
%
% \end{operationcontract}
%
% ---------------------------------------------------------------------------
% Referencing
% ---------------------------------------------------------------------------
%
% Requirements:
%
% \ref{fr:user-authentication}
% \ref{nfr:authentication-response-time}
%
% Use cases:
%
% \ref{uc:perform-login}
%
% Operation contracts:
%
% \ref{oc:create-order}
%
% The identifiers (FR01, NFR01, UC01, OC01, ...) are generated
% automatically, and labels are generated from the item names.
%
% ===========================================================================


% ===========================================================================
% PACKAGES
% ===========================================================================
\usepackage[utf8]{inputenc}
\usepackage[a4paper,margin=2cm]{geometry}
\usepackage{enumitem}
\usepackage{titlesec}
\usepackage{hyperref}
\usepackage{titling}
\usepackage{graphicx}
\usepackage[most]{tcolorbox}

% ---------------------------------------------------------------------------
% Customizations
% ---------------------------------------------------------------------------

% Subsection formatting
\titleformat{\subsection}
  {\normalfont\large\bfseries}
  {\thesubsection}
  {1em}
  {\underline}

% Itemize formatting
\setlist[itemize,1]{label=$\bullet$}
\setlist[itemize,2]{label=$\circ$}
\setlist[itemize,3]{label=$\diamond$}

% ---------------------------------------------------------------------------
% Title
% ---------------------------------------------------------------------------

\setlength{\droptitle}{-2cm}

\pretitle{\begin{flushleft}\LARGE\bfseries}
\posttitle{\end{flushleft}}

\preauthor{\begin{flushleft}\large}
\postauthor{\end{flushleft}}

\predate{\begin{flushleft}\large}
\postdate{\end{flushleft} \vspace{-1cm}}

% ---------------------------------------------------------------------------
% Requirements
% ---------------------------------------------------------------------------

\newcounter{functionalrequirement}
\newcounter{nonfunctionalrequirement}

% Functional requirement identifier: FR01, FR02, ...
\renewcommand{\thefunctionalrequirement}{%
    FR\ifnum\value{functionalrequirement}<10 0\fi\arabic{functionalrequirement}%
}

% Non-functional requirement identifier: NFR01, NFR02, ...
\renewcommand{\thenonfunctionalrequirement}{%
    NFR\ifnum\value{nonfunctionalrequirement}<10 0\fi\arabic{nonfunctionalrequirement}%
}

\ExplSyntaxOn

% Convert requirement name into a stable label.
\cs_new_protected:Npn \requirement_slug:n #1
{
    % Convert to lowercase.
    \tl_set:Nx \l_tmpa_tl
    {
        \str_lowercase:n { #1 }
    }

    \str_replace_all:Nnn \l_tmpa_tl { á } { a }
    \str_replace_all:Nnn \l_tmpa_tl { à } { a }
    \str_replace_all:Nnn \l_tmpa_tl { â } { a }
    \str_replace_all:Nnn \l_tmpa_tl { ã } { a }
    \str_replace_all:Nnn \l_tmpa_tl { ä } { a }

    \str_replace_all:Nnn \l_tmpa_tl { é } { e }
    \str_replace_all:Nnn \l_tmpa_tl { è } { e }
    \str_replace_all:Nnn \l_tmpa_tl { ê } { e }
    \str_replace_all:Nnn \l_tmpa_tl { ë } { e }

    \str_replace_all:Nnn \l_tmpa_tl { í } { i }
    \str_replace_all:Nnn \l_tmpa_tl { ì } { i }
    \str_replace_all:Nnn \l_tmpa_tl { î } { i }
    \str_replace_all:Nnn \l_tmpa_tl { ï } { i }

    \str_replace_all:Nnn \l_tmpa_tl { ó } { o }
    \str_replace_all:Nnn \l_tmpa_tl { ò } { o }
    \str_replace_all:Nnn \l_tmpa_tl { ô } { o }
    \str_replace_all:Nnn \l_tmpa_tl { õ } { o }
    \str_replace_all:Nnn \l_tmpa_tl { ö } { o }

    \str_replace_all:Nnn \l_tmpa_tl { ú } { u }
    \str_replace_all:Nnn \l_tmpa_tl { ù } { u }
    \str_replace_all:Nnn \l_tmpa_tl { û } { u }
    \str_replace_all:Nnn \l_tmpa_tl { ü } { u }

    \str_replace_all:Nnn \l_tmpa_tl { ç } { c }

    % Replace every sequence of non-alphanumeric characters with '-'.
    \regex_replace_all:nnN
        { [^a-z0-9]+ }
        { - }
        \l_tmpa_tl

    % Remove '-' from the beginning.
    \regex_replace_once:nnN
        { \A-+ }
        { }
        \l_tmpa_tl

    % Remove '-' from the end.
    \regex_replace_once:nnN
        { -+\Z }
        { }
        \l_tmpa_tl
}

% Functional requirement
\NewDocumentCommand{\fr}{m +m}
{
    \refstepcounter{functionalrequirement}
    \noindent
    \textbf{\thefunctionalrequirement{} - #1}
    \requirement_slug:n { #1 }
    \exp_args:Nx \label { fr:\l_tmpa_tl }

    \begin{itemize}[noitemsep, topsep=2pt, leftmargin=2em]
        #2
    \end{itemize}

    \par\medskip
}

% Non-functional requirement
\NewDocumentCommand{\nfr}{m +m}
{
    \refstepcounter{nonfunctionalrequirement}
    \noindent
    \textbf{\thenonfunctionalrequirement{} - #1}
    \requirement_slug:n { #1 }
    \exp_args:Nx \label { nfr:\l_tmpa_tl }

    \begin{itemize}[noitemsep, topsep=2pt, leftmargin=2em]
        #2
    \end{itemize}

    \par\medskip
}

\ExplSyntaxOff

% ---------------------------------------------------------------------------
% Use Cases
% ---------------------------------------------------------------------------

\newcounter{usecase}

% Use case identifier: UC01, UC02, ...
\renewcommand{\theusecase}{%
    UC\ifnum\value{usecase}<10 0\fi\arabic{usecase}%
}

% Use case box
\newtcolorbox{usecasebox}{
    enhanced,
    breakable,
    lines before break=15,
    colback=white,
    colframe=black!40,
    boxrule=0.5pt,
    arc=2pt,
    left=8pt,
    right=8pt,
    top=6pt,
    bottom=6pt,
    before skip=8pt,
    after skip=12pt,
    segmentation style={
        line width=0.5pt,
        solid
    }
}

\ExplSyntaxOn

% Use case environment
\NewDocumentEnvironment{usecase}{m}
{
    \refstepcounter{usecase}

    % Generate stable label from the use case name.
    \requirement_slug:n { #1 }
    \exp_args:Nx \label { uc:\l_tmpa_tl }

    % Compact title, similar to requirements.
    \begin{usecasebox}
    \noindent
    \textbf{\theusecase{} - #1}
    \par\smallskip
    \tcblower
}
{
    \end{usecasebox}
}

\ExplSyntaxOff

% Use Case Fields ----------------------------------------
\newcommand{\ucdescription}[1]{
    \noindent
    \textbf{Description}
    \par
    \leftskip=1em
    #1
    \par
    \leftskip=0pt
    \medskip
}

\newcommand{\ucactors}[1]{
    \noindent
    \textbf{Actors}
    \begin{itemize}[noitemsep, topsep=2pt, leftmargin=2em]
        #1
    \end{itemize}
}

\newcommand{\ucrelatedrequirements}[1]{
    \noindent
    \textbf{Related requirements}
    \begin{itemize}[noitemsep, topsep=2pt, leftmargin=2em]
        #1
    \end{itemize}
}

\newcommand{\ucpreconditions}[1]{
    \noindent
    \textbf{Preconditions}
    \begin{itemize}[noitemsep, topsep=2pt, leftmargin=2em]
        #1
    \end{itemize}
}

\newcommand{\ucmainflow}[1]{
    \noindent
    \textbf{Main Flow}
    \begin{enumerate}[noitemsep, topsep=2pt, leftmargin=2.4em]
        #1
    \end{enumerate}
}

\newcommand{\ucalternativeflow}[1]{
    \noindent
    \textbf{Alternative Flows}
    \begin{itemize}[noitemsep, topsep=2pt, leftmargin=2em]
        #1
    \end{itemize}
}

\newcommand{\ucbusinessrules}[1]{
    \noindent
    \textbf{Business Rules}
    \begin{itemize}[noitemsep, topsep=2pt, leftmargin=2em]
        #1
    \end{itemize}
}

\newcommand{\ucpostconditions}[1]{
    \noindent
    \textbf{Postconditions}
    \begin{itemize}[noitemsep, topsep=2pt, leftmargin=2em]
        #1
    \end{itemize}
}

% ---------------------------------------------------------------------------
% Operation Contracts
% ---------------------------------------------------------------------------

\newcounter{operationcontract}

% Operation contract identifier: OC01, OC02, ...
\renewcommand{\theoperationcontract}{%
    OC\ifnum\value{operationcontract}<10 0\fi\arabic{operationcontract}%
}

% Operation contract box
\newtcolorbox{operationcontractbox}{
    enhanced,
    breakable,
    lines before break=15,
    colback=white,
    colframe=black!40,
    boxrule=0.5pt,
    arc=2pt,
    left=8pt,
    right=8pt,
    top=6pt,
    bottom=6pt,
    before skip=8pt,
    after skip=12pt,
    segmentation style={
        line width=0.3pt,
        draw=black!30
    }
}

\ExplSyntaxOn

% Operation contract environment
\NewDocumentEnvironment{operationcontract}{m}
{
    \refstepcounter{operationcontract}

    % Generate stable label from the operation contract name.
    \requirement_slug:n { #1 }
    \exp_args:Nx \label { oc:\l_tmpa_tl }

    \begin{operationcontractbox}
        \noindent
        \textbf{\theoperationcontract{} - #1}

        \tcblower
}
{
    \end{operationcontractbox}
}

\ExplSyntaxOff

% Operation Contract Fields ------------------------------------------
\newcommand{\ocoperation}[1]{
    \noindent
    \textbf{Operation}
    \par
    \leftskip=1em
    \texttt{#1}
    \par
    \leftskip=0pt
    \medskip
}

\newcommand{\ocdescription}[1]{
    \noindent
    \textbf{Description}
    \par
    \leftskip=1em
    #1
    \par
    \leftskip=0pt
    \medskip
}

\newcommand{\ocpreconditions}[1]{
    \noindent
    \textbf{Preconditions}
    \begin{itemize}[noitemsep, topsep=2pt, leftmargin=2em]
        #1
    \end{itemize}
}

\newcommand{\ocpostconditions}[1]{
    \noindent
    \textbf{Postconditions}
    \begin{itemize}[noitemsep, topsep=2pt, leftmargin=2em]
        #1
    \end{itemize}
}

\newcommand{\ocrelatedusecases}[1]{
    \noindent
    \textbf{Related Use Cases}
    \begin{itemize}[noitemsep, topsep=2pt, leftmargin=2em]
        #1
    \end{itemize}
}

\usepackage[T1]{fontenc}
\usepackage[brazil]{babel}
\setlength{\parindent}{0pt}
\setlength{\parskip}{3pt plus 1pt minus 1pt}
\hypersetup{hidelinks,pdftitle={TP2: Refatoracao de um Sistema de Geracao de Relatorios},pdflang={pt-BR}}
\setlength{\emergencystretch}{4em}

\title{TP2: Refatoração de um Sistema de Geração de Relatórios\\[5pt]
\large Análise comparativa da implementação original e da arquitetura refatorada}
\author{}
\date{}

\begin{document}
\maketitle
\tableofcontents
\vspace{5pt}
\small

\section{Introdução}
O TP2 refatora um gerador de relatórios de cidades. O objetivo é preservar as saídas HTML e TXT e permitir novas fontes e formatos. A análise compara o código original com as alterações realizadas, destacando a redução de duplicação, a separação de responsabilidades, os princípios SOLID e os padrões de projeto.

A função do sistema permanece a mesma: ler cidades e gerar um relatório. A mudança está na arquitetura, que passa a separar \textbf{leitura, interpretação, formatação e coordenação}.

\section{Diagnóstico da implementação original}
No sistema original, \texttt{CidadesHTMLReport} e \texttt{CidadesTXTReporter} repetiam três etapas: \texttt{ler(caminho)} acessava o arquivo com \texttt{fs.readFileSync()}, \texttt{parse()} aplicava \texttt{JSON.parse()} e \texttt{reportar()} produzia a saída. Os dois primeiros métodos tinham a mesma implementação nas duas classes.

O atributo \texttt{this.cidades} armazenava inicialmente o texto lido e depois os objetos interpretados. Cada reporter acumulava acesso ao arquivo, interpretação e apresentação. Isso gerava \textbf{duplicação, baixo isolamento de responsabilidades e alto acoplamento}.

O \texttt{index.js} também repetia a criação e execução dos reporters em blocos condicionais. O caminho \texttt{./data/cidades-2.json} era fixo. Para adicionar uma saída ou fonte, seria necessário alterar várias partes do fluxo.

A refatoração substitui os dois fluxos completos \texttt{ler} $\rightarrow$ \texttt{parse} $\rightarrow$ \texttt{reportar} por componentes independentes de leitura e formatação.

\section{Refatorações realizadas}


\subsection{Extração da leitura para uma abstração comum}
A leitura comum foi extraída para \texttt{ILeitor}. Seu método \texttt{ler(caminho)} lê o arquivo e chama \texttt{this.parse(conteudo)}. Assim, a operação compartilhada fica em um único lugar.

\texttt{ILeitor} funciona como classe abstrata. As subclasses implementam \texttt{parse()} conforme a fonte. Os formatadores deixam de acessar arquivos e interpretar JSON. \textbf{A duplicação da leitura foi eliminada.}

\subsection{Especialização das fontes de dados}
\texttt{LeitorJSON} utiliza \texttt{JSON.parse()} para interpretar \texttt{cidades-2.json}. Já \texttt{LeitorCSV} processa as linhas de \texttt{cidades.csv} e cria objetos com os campos \texttt{Estado}, \texttt{ID} e \texttt{Nome}.

Ambos entregam dados no mesmo formato interno. Por isso, qualquer fonte pode ser combinada com HTML, TXT ou CSV. \textbf{A origem dos dados não define mais o formato do relatório.}

\subsection{Extração da formatação para classes independentes}
O método \texttt{reportar()} foi substituído por formatadores independentes: \texttt{FormatadorHTML} e \texttt{FormatadorTXT}. Ambos estendem \texttt{IFormatador} e implementam \texttt{formatar(dados)}. A classe base define o contrato de saída, sem assumir a leitura dos dados.

Também foi criado \texttt{FormatadorCSV}. Ele gera as colunas \texttt{Estado}, \texttt{ID} e \texttt{Nome} e trata aspas, vírgulas e quebras de linha nos campos. O novo formato foi incluído \textbf{sem modificar HTML, TXT ou o gerador}.

\subsection{Introdução de uma classe coordenadora}
Foi criada \texttt{GeradorDeRelatorio}, que recebe leitor e formatador no construtor. Em \texttt{gerar(caminho)}, chama \texttt{leitor.ler(caminho)} e encaminha o resultado a \texttt{formatador.formatar(dados)}.

Antes, cada reporter executava todas as etapas. Agora, os leitores obtêm os dados, os formatadores produzem a saída e o gerador coordena a execução. \textbf{A composição reduz o acoplamento e evita repetir o fluxo.}

\subsection{Seleção de fonte e formato pela linha de comando}
O \texttt{index.js} original selecionava apenas HTML ou TXT e utilizava um arquivo fixo. Na nova versão, recebe dois argumentos de \texttt{process.argv}: caminho do arquivo e formato desejado.

A extensão seleciona \texttt{LeitorJSON} ou \texttt{LeitorCSV}. O segundo argumento seleciona \texttt{FormatadorHTML}, \texttt{FormatadorTXT} ou \texttt{FormatadorCSV}. Os registros de classes substituem os blocos repetidos. O programa cria um único \texttt{GeradorDeRelatorio} com os componentes escolhidos.

Exemplos: \texttt{node index.js ./data/cidades.csv html} e \texttt{node index.js ./data/cidades-2.json csv}. \textbf{O ponto de entrada configura a execução; não lê nem formata os dados.}

\subsection{Correção pontual do relatório TXT}
O relatório TXT apresentava a primeira cidade junto à linha de sinais de igualdade do cabeçalho. A correção proposta adiciona \verb|\n| após o cabeçalho de \texttt{FormatadorTXT.js}.

A mudança fica restrita ao formatador TXT, sem afetar os leitores ou os demais formatos.

\subsection{Tratamento de entradas inválidas e falhas de execução}
O \texttt{index.js} passou a validar a extensão do arquivo e o formato solicitado. Se uma opção for inválida, informa o uso correto e encerra com erro.

A geração também foi envolvida por \texttt{try/catch}. O erro \texttt{ENOENT} recebe uma mensagem de arquivo não encontrado; outras falhas exibem a descrição da exceção. As mensagens utilizam \texttt{console.error()}, preservando a saída padrão para o relatório, inclusive quando redirecionada com \texttt{>}.

\textbf{O tratamento centralizado facilita o diagnóstico sem duplicar verificações nos formatadores.}

\section{Arquitetura refatorada e padrões de projeto}
A arquitetura original possuía dois reporters completos, cada um responsável por leitura, interpretação e apresentação. A nova estrutura divide essas tarefas entre \texttt{ILeitor} e suas subclasses, \texttt{IFormatador} e suas subclasses, e \texttt{GeradorDeRelatorio}.

O sistema passa a combinar componentes independentes, em vez de criar uma classe para cada par de fonte e saída.

\subsection{Template Method na leitura}
O padrão \textbf{Template Method} aparece em \texttt{ILeitor}. O método \texttt{ler()} fixa a sequência: acessar o arquivo e interpretar seu conteúdo. A etapa variável, \texttt{parse()}, é implementada por \texttt{LeitorJSON} e \texttt{LeitorCSV}.

A leitura comum é reutilizada, enquanto cada fonte define seu próprio parsing. \texttt{IFormatador} não representa esse padrão: apenas estabelece o contrato \texttt{formatar()}.

\subsection{Strategy na escolha dos componentes}
O padrão \textbf{Strategy} permite trocar leitores e formatadores sem alterar \texttt{GeradorDeRelatorio}. Por exemplo, \texttt{LeitorJSON} pode ser substituído por \texttt{LeitorCSV}; \texttt{FormatadorHTML}, por \texttt{FormatadorTXT} ou \texttt{FormatadorCSV}.

O gerador utiliza sempre os métodos \texttt{ler()} e \texttt{formatar()}. As estratégias são escolhidas na entrada da aplicação e injetadas no construtor. \textbf{Entrada e saída podem variar de forma independente.}

\subsection{Delimitação dos padrões reivindicados}
A separação entre leitura e formatação lembra o \textbf{Bridge}, mas não caracteriza, por si só, sua aplicação canônica. Os padrões efetivamente demonstrados são \textbf{Template Method} e \textbf{Strategy}.

Não há uma classe Factory implementada. Sua inclusão é opcional e não é necessária para atender ao objetivo da refatoração.

\section{Relação com os princípios SOLID}


\subsection{SRP --- Single Responsibility Principle}
O \textbf{SRP} é atendido pela divisão das responsabilidades. Antes, os reporters liam, interpretavam e formatavam. Agora, \texttt{ILeitor} centraliza a leitura; as subclasses interpretam os dados; os formatadores produzem a saída; e \texttt{GeradorDeRelatorio} coordena as etapas.

Cada componente passa a ter um motivo principal para mudar.

\subsection{OCP --- Open/Closed Principle}
O \textbf{OCP} é favorecido pela extensão com novas classes. Um \texttt{FormatadorMarkdown}, por exemplo, pode implementar \texttt{formatar(dados)} sem modificar o gerador ou os formatadores existentes. Novas fontes seguem o mesmo modelo por meio de \texttt{ILeitor}.

Ainda é necessário registrar a nova classe no \texttt{index.js}. Portanto, a abertura para extensão é maior no núcleo do sistema do que na configuração inicial.

\subsection{LSP --- Liskov Substitution Principle}
O \textbf{LSP} é observado quando \texttt{LeitorJSON} e \texttt{LeitorCSV} podem ocupar o lugar de \texttt{ILeitor}, retornando coleções compatíveis. O mesmo vale para os formatadores, que recebem os dados e retornam uma string.

A substituição depende do contrato: \texttt{cidades-1.json}, por exemplo, possui estrutura diferente e não é adaptado pelo leitor JSON atual.

\subsection{ISP --- Interface Segregation Principle}
O \textbf{ISP} é refletido em contratos pequenos. Leitores não precisam implementar formatação; formatadores não precisam implementar leitura.

Em JavaScript, \texttt{ILeitor} e \texttt{IFormatador} são classes abstratas, não interfaces formais. Ainda assim, promovem a separação das responsabilidades.

\subsection{DIP --- Dependency Inversion Principle}
O \textbf{DIP} aparece na injeção de dependências por construtor. \texttt{GeradorDeRelatorio} recebe leitor e formatador, em vez de criar classes concretas. Sua lógica depende dos métodos \texttt{ler()} e \texttt{formatar()}.

A compatibilidade é verificada dinamicamente. Imports não utilizados de \texttt{ILeitor} e \texttt{IFormatador} no gerador podem ser removidos.

\textbf{Síntese:} SRP e DIP são as melhorias mais evidentes. OCP, LSP e ISP são favorecidos pela separação e substituição dos componentes.

\section{Verificação funcional e preservação de comportamento}
A finalidade do programa foi preservada. Os componentes foram testados com \texttt{cidades-2.json} (5.340 registros) e \texttt{cidades.csv} (5.570 registros). Durante o desenvolvimento, foi confirmado o funcionamento das \textbf{seis combinações} de entrada JSON/CSV e saída HTML/TXT/CSV.

A execução utiliza o mesmo \texttt{GeradorDeRelatorio} em todos os casos. Como as fontes contêm registros diferentes, a validação considera os dados de cada arquivo, não uma igualdade entre as bases.

Também devem ser verificados os casos de formato inválido, arquivo inexistente e a primeira linha do TXT. Esses testes podem ser automatizados antes da entrega.

\section{Limitações e estado da implementação}
A arquitetura está representada por \texttt{ILeitor}, \texttt{LeitorJSON}, \texttt{LeitorCSV}, \texttt{IFormatador}, \texttt{FormatadorHTML}, \texttt{FormatadorTXT}, \texttt{FormatadorCSV} e \texttt{GeradorDeRelatorio}. A seleção por linha de comando foi confirmada nas seis combinações.

A correção do TXT e o tratamento de erros foram propostos em código, mas os arquivos finais não foram reenviados para conferência. Sua integração ao repositório deve ser verificada.

O \texttt{LeitorCSV} utiliza \texttt{split(',')}, adequado aos dados fornecidos, mas insuficiente para campos CSV com vírgulas entre aspas. O \texttt{cidades-1.json} também exige adaptação específica. Essas limitações não impedem o atendimento aos requisitos básicos.

\section{Conclusão}
A refatoração eliminou a duplicação da leitura e separou obtenção, interpretação e apresentação dos dados. O \texttt{GeradorDeRelatorio} passou a coordenar essas etapas por composição. O \texttt{index.js} passou a selecionar fonte e formato por argumentos.

Os padrões \textbf{Template Method} e \textbf{Strategy} sustentam a nova estrutura. Os princípios SOLID aparecem principalmente na responsabilidade única, na injeção de dependências e na facilidade de extensão.

O suporte a CSV como entrada e saída, somado às seis combinações funcionais, demonstra o resultado da mudança. A correção do TXT e o tratamento de erros completam os ajustes de qualidade, sujeitos à verificação no repositório.

\emph{Base da análise: código JavaScript, diagramas Mermaid e arquivos JSON/CSV fornecidos para o TP2.}

\end{document}
